% Options for packages loaded elsewhere
\PassOptionsToPackage{unicode}{hyperref}
\PassOptionsToPackage{hyphens}{url}
\documentclass[
  a4paper]{article}
\usepackage{xcolor}
\usepackage[margin=2.5cm]{geometry}
\usepackage{amsmath,amssymb}
\setcounter{secnumdepth}{-\maxdimen} % remove section numbering
\usepackage{iftex}
\ifPDFTeX
  \usepackage[T1]{fontenc}
  \usepackage[utf8]{inputenc}
  \usepackage{textcomp} % provide euro and other symbols
\else % if luatex or xetex
  \usepackage{unicode-math} % this also loads fontspec
  \defaultfontfeatures{Scale=MatchLowercase}
  \defaultfontfeatures[\rmfamily]{Ligatures=TeX,Scale=1}
\fi
\usepackage{lmodern}
\ifPDFTeX\else
  % xetex/luatex font selection
\fi
% Use upquote if available, for straight quotes in verbatim environments
\IfFileExists{upquote.sty}{\usepackage{upquote}}{}
\IfFileExists{microtype.sty}{% use microtype if available
  \usepackage[]{microtype}
  \UseMicrotypeSet[protrusion]{basicmath} % disable protrusion for tt fonts
}{}
\makeatletter
\@ifundefined{KOMAClassName}{% if non-KOMA class
  \IfFileExists{parskip.sty}{%
    \usepackage{parskip}
  }{% else
    \setlength{\parindent}{0pt}
    \setlength{\parskip}{6pt plus 2pt minus 1pt}}
}{% if KOMA class
  \KOMAoptions{parskip=half}}
\makeatother
\usepackage{longtable,booktabs,array}
\newcounter{none} % for unnumbered tables
\usepackage{calc} % for calculating minipage widths
% Correct order of tables after \paragraph or \subparagraph
\usepackage{etoolbox}
\makeatletter
\patchcmd\longtable{\par}{\if@noskipsec\mbox{}\fi\par}{}{}
\makeatother
% Allow footnotes in longtable head/foot
\IfFileExists{footnotehyper.sty}{\usepackage{footnotehyper}}{\usepackage{footnote}}
\makesavenoteenv{longtable}
\setlength{\emergencystretch}{3em} % prevent overfull lines
\providecommand{\tightlist}{%
  \setlength{\itemsep}{0pt}\setlength{\parskip}{0pt}}
\usepackage{bookmark}
\IfFileExists{xurl.sty}{\usepackage{xurl}}{} % add URL line breaks if available
\urlstyle{same}
\hypersetup{
  hidelinks,
  pdfcreator={LaTeX via pandoc}}

\author{}
\date{}

\begin{document}

\section{\texorpdfstring{Regularity of Central Projection Between
Hyperspherical Shells: A Formulation Without Tangent Hyperplanes and
Degeneracy at
\(R = 0\)}{Regularity of Central Projection Between Hyperspherical Shells: A Formulation Without Tangent Hyperplanes and Degeneracy at R = 0}}\label{regularity-of-central-projection-between-hyperspherical-shells-a-formulation-without-tangent-hyperplanes-and-degeneracy-at-r-0}

\textbf{Author}: Noriaki Kihara\\
\textbf{Affiliation}: WF System Co., Ltd.\\
\textbf{ORCID}: 0009-0004-6753-4020\\
\textbf{Version}: v1.0\\
\textbf{Date}: April 2026\\
\textbf{DOI}:
\href{https://doi.org/10.5281/zenodo.19692192}{10.5281/zenodo.19692192}

\begin{center}\rule{0.5\linewidth}{0.5pt}\end{center}

\subsection{Abstract}\label{abstract}

In the classical formulation {[}P1{]}, central projection is defined via
a tangent hyperplane, where the divergence of \(\tan\theta\) at the
equator necessitates the constraint \(R > 0\). This paper proves that
central projection defined directly between surfaces of concentric
hyperspherical shells does not require this constraint.

Specifically, we show:

\begin{enumerate}
\def\labelenumi{\arabic{enumi}.}
\tightlist
\item
  The interspherical central projection
  \(\Phi_{R \to R'} : S^n(R) \to S^n(R')\) is a diffeomorphism for any
  \(R, R' > 0\), regular and non-divergent over the entire sphere
  (Theorem 1).
\item
  The projection to \(R' = 0\) is well-defined and smooth, with no
  divergence. The only limitation is that the Jacobian degenerates to
  zero and the inverse map becomes undefined (Theorem 2).
\item
  The classical gnomonic projection via a tangent hyperplane genuinely
  diverges at the equator. This divergence is intrinsic to the use of
  the tangent hyperplane and does not structurally exist in
  interspherical projection (Theorem 3).
\end{enumerate}

These results confirm that the constraint \(R, R_1, R_0 > 0\) in the
three-layer structure model {[}M1{]} is unnecessary. All claims of
{[}M1{]} hold over the entire real domain including \(R = 0\), where the
only limitation is the undefined inverse map (degeneracy).

\textbf{Keywords}: central projection, interspherical projection,
regularity, gnomonic projection, equatorial divergence, degeneracy,
diffeomorphism

\begin{center}\rule{0.5\linewidth}{0.5pt}\end{center}

\subsection{§1 Introduction and
Motivation}\label{introduction-and-motivation}

\subsubsection{§1.1 Limitations of Formulation via Tangent
Hyperplanes}\label{limitations-of-formulation-via-tangent-hyperplanes}

In {[}P1{]}, central projection was defined as the projection from the
origin \(O\) of the ambient coordinate space \(\mathbb{R}^{n+1}\) onto
the hypersphere \(S^n(R)\):

\[
P_R : \mathbb{R}^{n+1} \setminus \{O\} \to S^n(R), \quad P_R(\mathbf{x}) = R \cdot \frac{\mathbf{x}}{\|\mathbf{x}\|}
\]

This definition is meaningful only for \(R > 0\). The requirement
\(R > 0\) arises for two reasons:

\begin{enumerate}
\def\labelenumi{\arabic{enumi}.}
\tightlist
\item
  \textbf{Formal reason}: When \(R = 0\), \(S^n(0) = \{O\}\) degenerates
  and the image becomes a single point.
\item
  \textbf{Pathological properties of tangent hyperplanes}: In the
  formulation of {[}P1{]}, the pullback metric is constructed via the
  tangent hyperplane \(T\) at the north pole of \(S^n(R)\). The inverse
  map through this tangent hyperplane diverges as
  \(\tan\theta \to \infty\) at the equator (\(\theta = \pi/2\)) and does
  not reach the southern hemisphere.
\end{enumerate}

This paper proves that the second reason is \textbf{a limitation arising
from the use of the tangent hyperplane, which does not exist in
interspherical central projection}, and that the first reason is merely
a limited restriction of ``undefined inverse map.''

\subsubsection{§1.2 Observation on the Three-Layer Structure
Model}\label{observation-on-the-three-layer-structure-model}

{[}M1{]} defined the central projection \(\Phi_{1 \to 0}\) between
hyperspheres \(S^n(R_1)\) and \(S^n(R_0)\) of different curvature radii
and asserted that this map is well-defined and regular. All claims of
{[}M1{]} use only interspherical projection and never employ tangent
hyperplanes.

This observation raises a natural question: Is the constraint \(R > 0\)
in {[}M1{]} essential, or merely a remnant from the formulation of
{[}P1{]}?

This paper proves that the constraint is merely a remnant, and that
interspherical central projection is regular over the entire real domain
including \(R = 0\).

\begin{center}\rule{0.5\linewidth}{0.5pt}\end{center}

\subsection{§2 Definition of Central Projection Between Hyperspherical
Shells}\label{definition-of-central-projection-between-hyperspherical-shells}

\subsubsection{Definition 2.1 (Interspherical Central
Projection)}\label{definition-2.1-interspherical-central-projection}

For \(R, R' > 0\), the central projection between concentric
hyperspheres \(S^n(R)\) and \(S^n(R')\) is defined as:

\[
\Phi_{R \to R'} : S^n(R) \to S^n(R'), \quad \Phi_{R \to R'}(p) = \frac{R'}{R} \cdot p
\]

where \(p \in S^n(R) \subset \mathbb{R}^{n+1}\) with \(\|p\| = R\).

\textbf{Remark}: This map scales \(p\) along the half-line through \(O\)
to distance \(R'\). It is equivalent to
\(\Phi_{1 \to 0} = P_{R_0} \circ P_{R_1}^{-1}\) from {[}M1{]} §5.3, but
here we use the direct representation.

\subsubsection{\texorpdfstring{Definition 2.2 (Extension to
\(R' = 0\))}{Definition 2.2 (Extension to R\textquotesingle{} = 0)}}\label{definition-2.2-extension-to-r-0}

For \(R > 0\), \(R' = 0\), we define:

\[
\Phi_{R \to 0} : S^n(R) \to \{O\}, \quad \Phi_{R \to 0}(p) = O \quad (\forall p \in S^n(R))
\]

This is the constant map naturally obtained as the limit \(R' \to 0\) in
Definition 2.1.

\subsubsection{Definition 2.3
(Degeneracy)}\label{definition-2.3-degeneracy}

A map \(\Phi : X \to Y\) is said to be \textbf{degenerate} at
\(x_0 \in X\) if \(\Phi\) is smooth in a neighborhood of \(x_0\) but the
Jacobian vanishes at \(x_0\). Degeneracy does not signify loss of
smoothness (singularity) but rather the local undefined nature of the
inverse map.

\begin{center}\rule{0.5\linewidth}{0.5pt}\end{center}

\subsection{§3 Main Results}\label{main-results}

\subsubsection{Theorem 1 (Diffeomorphism of Interspherical
Projection)}\label{theorem-1-diffeomorphism-of-interspherical-projection}

For \(R, R' > 0\), \(\Phi_{R \to R'} : S^n(R) \to S^n(R')\) is a
diffeomorphism. In particular:

\begin{enumerate}
\def\labelenumi{(\roman{enumi})}
\tightlist
\item
  \(\Phi_{R \to R'}\) is well-defined, smooth, and bijective over the
  entire sphere.\\
\item
  The inverse map \(\Phi_{R \to R'}^{-1} = \Phi_{R' \to R}\) is also
  smooth.\\
\item
  The Jacobian is the constant \((R'/R)^n\) over the entire sphere and
  never vanishes.\\
\item
  The image is all of \(S^n(R')\); no hemisphere restriction exists.
\end{enumerate}

\textbf{Proof}:

Represent points on \(S^n(R)\) as \(p = R \cdot \hat{u}\)
(\(\hat{u} \in S^n(1)\), \(\|\hat{u}\| = 1\)).

\begin{enumerate}
\def\labelenumi{(\roman{enumi})}
\item
  \(\Phi_{R \to R'}(p) = (R'/R) \cdot p = R' \cdot \hat{u}\). Since
  \(\|R' \cdot \hat{u}\| = R'\), the image lies on \(S^n(R')\). The
  image is unique for each \(\hat{u}\), hence well-defined; being a
  scalar multiple, it is smooth.
  \(\hat{u} \neq \hat{u}' \Rightarrow R' \hat{u} \neq R' \hat{u}'\)
  gives injectivity; for any \(q = R' \hat{v} \in S^n(R')\),
  \(p = R \hat{v}\) provides the preimage, giving surjectivity.
\item
  \(\Phi_{R' \to R}(q) = (R/R') \cdot q\), and
  \(\Phi_{R' \to R} \circ \Phi_{R \to R'} = \mathrm{id}_{S^n(R)}\) is
  verified by direct computation. Since \(R/R'\) is a finite nonzero
  constant, the inverse map is also smooth.
\item
  Since \(S^n(R)\) is the \(R\)-scaled copy of \(S^n(1)\), using local
  coordinates \((\theta_1, \ldots, \theta_n)\) on the unit sphere
  \(S^n(1)\), the induced metric on \(S^n(R)\) is
  \(g_{ij}^{(R)} = R^2 \cdot g_{ij}^{(1)}\). Since \(\Phi_{R \to R'}\)
  induces the identity map \(\hat{u} \mapsto \hat{u}\) on the unit
  sphere, the Jacobian matrix of \(\Phi_{R \to R'}\) gives a scaling of
  \(R'/R\) in each tangent direction. The \(n\)-dimensional Jacobian
  (ratio of volume elements) is \((R'/R)^n\), which never vanishes for
  \(R, R' > 0\).
\item
  Directly from surjectivity in (i). Since \(\hat{u}\) ranges over the
  entire unit sphere, the image covers all of \(S^n(R')\).
  \(\blacksquare\)
\end{enumerate}

\begin{center}\rule{0.5\linewidth}{0.5pt}\end{center}

\subsubsection{\texorpdfstring{Theorem 2 (Regularity and Degeneracy at
\(R' = 0\))}{Theorem 2 (Regularity and Degeneracy at R\textquotesingle{} = 0)}}\label{theorem-2-regularity-and-degeneracy-at-r-0}

For \(R > 0\), the following hold for
\(\Phi_{R \to 0} : S^n(R) \to \{O\}\):

\begin{enumerate}
\def\labelenumi{(\roman{enumi})}
\tightlist
\item
  \(\Phi_{R \to 0}\) is well-defined and smooth.\\
\item
  The image is bounded over the entire sphere
  (\(\|\Phi_{R \to 0}(p)\| = 0\), \(\forall p\)), and no
  \(\tan\theta\)-type divergence as in gnomonic projection exists.\\
\item
  The Jacobian degenerates to \(0\) over the entire sphere.\\
\item
  \(\Phi_{R \to 0}\) is not bijective, and the inverse map
  \(\Phi_{0 \to R}\) is undefined.
\end{enumerate}

\textbf{Proof}:

\begin{enumerate}
\def\labelenumi{(\roman{enumi})}
\item
  \(\Phi_{R \to 0}(p) = (0/R) \cdot p = \mathbf{0} = O\) is a constant
  map, clearly well-defined and smooth (\(C^\infty\)).
\item
  For any \(p \in S^n(R)\), \(\|\Phi_{R \to 0}(p)\| = 0 < \infty\), so
  the image is bounded. The unbounded behavior
  \(d(\theta) = R\tan\theta \to \infty\) of gnomonic projection (Theorem
  3) does not structurally exist.
\item
  In the expression from Theorem 1 (iii), the Jacobian is \((R'/R)^n\).
  For \(R' = 0\), \((0/R)^n = 0\). The degeneracy of the Jacobian to
  zero reflects that the map collapses all points to a single point; it
  is not a divergence.
\item
  \(\Phi_{R \to 0}\) maps all points of \(S^n(R)\) to \(O\), so it is
  not injective. Therefore, the inverse map does not exist.
  Specifically, the preimage of \(O\) is the entirety of \(S^n(R)\),
  making it impossible to determine which \(p \in S^n(R)\) should be
  recovered from \(O\). \(\blacksquare\)
\end{enumerate}

\begin{center}\rule{0.5\linewidth}{0.5pt}\end{center}

\subsubsection{Theorem 3 (Comparison with Projection via Tangent
Hyperplane)}\label{theorem-3-comparison-with-projection-via-tangent-hyperplane}

The inverse map \(G^{-1} : S^n(R) \to T_N\) of the classical gnomonic
projection

\[
G : T_N \to S^n(R), \quad G(\mathbf{y}) = R \cdot \frac{\mathbf{y}}{\|\mathbf{y}\|}
\]

via the tangent hyperplane \(T_N\) at the north pole
\(N = (0, \ldots, 0, R)\) of \(S^n(R)\), diverges at the equator (zenith
angle \(\theta = \pi/2\)).

\textbf{Proof}:

Represent points on \(S^n(R)\) by the zenith angle \(\theta\) (angular
distance from the north pole). The distance from the north pole to the
corresponding point on the tangent hyperplane \(T_N\) is

\[
d(\theta) = R \tan\theta
\]

As \(\theta \to \pi/2\), \(\tan\theta \to +\infty\), so
\(d(\theta) \to +\infty\). Therefore, \(G^{-1}\) diverges at the
equator.

Furthermore, for \(\theta > \pi/2\) (southern hemisphere), \(G^{-1}\) is
undefined. Since the tangent hyperplane \(T_N\) lies only in the
half-space on the north pole side, no preimage on \(T_N\) exists for
points on the southern hemisphere. Even extending the tangent hyperplane
to both sides does not resolve the equatorial divergence, nor does it
achieve coverage of the southern hemisphere. This is because the
divergence is intrinsically due to the geometry of the tangent
hyperplane (its placement as a tangent space of the hypersphere).

\textbf{Comparison}: The interspherical projection \(\Phi_{R \to R'}\)
(Theorem 1) does not structurally possess this divergence.
\(\Phi_{R \to R'}\) is a scalar multiplication and does not involve
unbounded functions such as \(\tan\theta\). The Jacobian is uniformly
\((R'/R)^n\) over the entire sphere and has no dependence on \(\theta\).
\(\blacksquare\)

\begin{center}\rule{0.5\linewidth}{0.5pt}\end{center}

\subsection{§4 Consequences for Prior
Works}\label{consequences-for-prior-works}

\subsubsection{\texorpdfstring{§4.1 Removal of the \(R > 0\) Constraint
in the Three-Layer Structure Model
{[}M1{]}}{§4.1 Removal of the R \textgreater{} 0 Constraint in the Three-Layer Structure Model {[}M1{]}}}\label{removal-of-the-r-0-constraint-in-the-three-layer-structure-model-m1}

In {[}M1{]}, the condition \(R > 0\) was imposed on the three curvature
radii \(R\), \(R_1\), \(R_0\) ({[}M1{]} §2 Condition 1, §3.1). This
constraint was inherited from the formulation of {[}P1{]} (\(R > 0\)).

By Theorems 1 and 2 of this paper, the claims of {[}M1{]} are extended
as follows:

\textbf{Corollary 4.1}: Claims 4.1 (shared origin), 4.2 (collinearity of
projection axes), 4.3 (concentric nested structure), and 5.1 (metric
non-contribution of \(R\)-directional displacement) of {[}M1{]} hold
over the entire domain \(R, R_1, R_0 \geq 0\). The only limitation is
that the inverse map \(\Phi_{0 \to R}\) becomes undefined for
projections involving \(R = 0\).

\textbf{Proof}: Each claim of {[}M1{]} relies solely on properties of
interspherical central projection and does not use tangent hyperplanes.
By Theorem 1, claims for \(R, R' > 0\) hold as stated. For \(R' = 0\),
Theorem 2 (i)(ii) shows that the forward projection is well-defined and
regular, so the claims of shared origin, collinearity, and concentric
structure are unaffected. Metric non-contribution (Claim 5.1) holds
trivially since the tangent space of \(S^n(0) = \{O\}\) degenerates to
zero dimension, making the orthogonality assertion vacuously true.
\(\blacksquare\)

\textbf{Remark}: This corollary is a mathematical extension of scope and
does not override the non-claim explicitly stated in {[}M1{]} §8.1 that
``we do not assert the picture of particles condensing at the
\(R \to 0\) limit.'' The regularity at \(R = 0\) means only that the map
is mathematically well-defined; it does not assert physical condensation
phenomena. Even in the cases \(R_1 = 0\) (particles condensing to the
origin) or \(R_0 = 0\) (our spacetime condensing to the origin), the
forward projection is mathematically well-defined as a constant map, but
their physical interpretations lie outside the scope of this paper and
{[}M1{]}.

\subsubsection{§4.2 Relationship with the Formulation of
{[}P1{]}}\label{relationship-with-the-formulation-of-p1}

The formulation via tangent hyperplanes in {[}P1{]} was essential for
constructing the pullback metric and deriving the Einstein tensor. This
paper does not claim that the formulation of {[}P1{]} is erroneous.

What this paper claims is that \textbf{when using only interspherical
central projection} (the framework of {[}M1{]}), the \(R > 0\)
constraint of {[}P1{]} is unnecessary. The use of the tangent hyperplane
is required for the specific purposes of {[}P1{]} (metric construction)
and is unnecessary for the geometric structures of interspherical
projection (shared origin, nested structure, direction preservation).

\begin{center}\rule{0.5\linewidth}{0.5pt}\end{center}

\subsection{§5 Conclusion}\label{conclusion}

This paper has proved that central projection defined between surfaces
of concentric hyperspherical shells possesses none of the pathological
properties of the classical gnomonic projection (projection via tangent
hyperplanes).

{\def\LTcaptype{none} % do not increment counter
\begin{longtable}[]{@{}
  >{\raggedright\arraybackslash}p{(\linewidth - 4\tabcolsep) * \real{0.1205}}
  >{\raggedright\arraybackslash}p{(\linewidth - 4\tabcolsep) * \real{0.4096}}
  >{\raggedright\arraybackslash}p{(\linewidth - 4\tabcolsep) * \real{0.4699}}@{}}
\toprule\noalign{}
\begin{minipage}[b]{\linewidth}\raggedright
Property
\end{minipage} & \begin{minipage}[b]{\linewidth}\raggedright
Via tangent hyperplane (gnomonic)
\end{minipage} & \begin{minipage}[b]{\linewidth}\raggedright
Interspherical projection (this paper)
\end{minipage} \\
\midrule\noalign{}
\endhead
\bottomrule\noalign{}
\endlastfoot
Regular domain & Northern hemisphere only (\(\theta < \pi/2\)) & Entire
sphere (\(\theta \in [0, \pi]\)) \\
Behavior at equator & Diverges as \(\tan\theta \to \infty\) & No
divergence \\
Jacobian & Depends on \(\theta\), diverges at equator & Constant
\((R'/R)^n\), uniform over sphere \\
\(R' = 0\) & Undefined as projection target (tangent hyperplane cannot
be constructed) & No divergence; only inverse map undefined \\
Bijectivity & Limited to northern hemisphere & \(R' > 0\): bijective
over entire sphere; \(R' = 0\): not bijective \\
\end{longtable}
}

Interspherical central projection is regular (non-divergent) over the
entire domain \(R \geq 0\). The only limitation at \(R = 0\) is the
undefined inverse map due to degeneracy of the Jacobian to zero.

These results confirm that the constraint \(R > 0\) in the three-layer
structure model {[}M1{]} is unnecessary. The geometric claims of
{[}M1{]} hold over the entire domain including \(R = 0\).

\begin{center}\rule{0.5\linewidth}{0.5pt}\end{center}

\subsection{References}\label{references}

\begin{itemize}
\tightlist
\item
  {[}P1{]}: Noriaki Kihara, Geometric Formulation of Four-Dimensional
  Spacetime via Central Projection, Zenodo, 2026. DOI:
  \href{https://doi.org/10.5281/zenodo.19427780}{10.5281/zenodo.19427780}
\item
  {[}M1{]}: Noriaki Kihara, Three-Layer Structure Model of Central
  Projection: Nested Structure of Subjective Spaces via
  \(R\)--\(R_1\)--\(R_0\) and Middleware Geometry, Zenodo, 2026. DOI:
  \href{https://doi.org/10.5281/zenodo.19691713}{10.5281/zenodo.19691713}
\end{itemize}

\begin{center}\rule{0.5\linewidth}{0.5pt}\end{center}

\emph{v1.0}

\end{document}
